% !TeX root = ../教材题库_全部条目_转写工作稿.tex
% Source: XJ-B1-C04-EX-P099-Y1526
% 第4章 / 4.1 实数指数幂和幂函数 / 教材 p.99、100 / PDF 第106、107页
% OCR未校对：公式与文字均可能有误；未完成数学TeX转写；禁止直接发给学生
% 原题及相邻原解截图（核对用）：../题目原图/XJ-B1-C04-EX-P099-Y1526-01.webp, ../题目原图/XJ-B1-C04-EX-P099-Y1526-02.webp
% 本文件只在转写完成后作为题干使用；OCR 条目尚不是可靠题干。
\begin{bankitem}{4.1 例题 例题7}{第4章 / 4.1 实数指数幂和幂函数 / 教材 p.99、100 / PDF 第106、107页}{XJ-B1-C04-EX-P099-Y1526}{OCR未校对：公式与文字均可能有误；未完成数学TeX转写；禁止直接发给学生}
\begin{ocrdraft}
[教材p.99 / PDF第106页]
\par
0才已知>1
\par
对任意的实数如求证："一2/'
\par
证明因为磊
\par
a，都是正数，且磊
\par
[教材p.100 / PDF第107页]
\par
故a
\par
“十h
\par
一也是正数．
\par
又
\par
u十2h
\par
(a
\par
一ah>1，
\par
即得
\end{ocrdraft}
\end{bankitem}
